\begin{lemma}
For every $q\in\mathbb N$, the number of independent pairs in the
neighborhood of $\ast$ in $L(A_q)$ is $9q^2$.
\end{lemma}
\begin{proof}
By \noderef{IndependentPairCount}, the objects counted are two-element
sets of neighbors whose two distinct indices are not adjacent.
By \noderef{OrderThreeAffineGeometry} the neighbors are exactly the
nonvertical indices. For distinct indices, \noderef{LineGraphOfHypergraph}
identifies nonadjacency with disjointness of their edge sets. The disjointness
clause of \noderef{OrderThreeAffineGeometry} therefore says that an
independent pair consists of indices with one common slope and two distinct
intercepts.

Here is an explicit bijection valid for arbitrary $q$. Let
$\mathcal J=\{(0,1),(0,2),(1,2)\}$, using the indicated representatives of
the residues. Map
\[
(m,(a,b),i,j)\longmapsto\{(m,a,i),(m,b,j)\}
\]
from $(\mathbb Z/3\mathbb Z)\times\mathcal J\times
\{0,\ldots,q-1\}^2$ to the independent pairs. Different intercepts
ensure that the image has two elements, and the geometry ensures independence.
Conversely, take any independent pair. Its common slope determines $m$;
sort its two intercepts using $0<1<2$ to obtain the unique $(a,b)\in\mathcal J$.
The unique edge-index of the pair above $a$ determines $i$, and the unique
one above $b$ determines $j$. This defines an inverse, since reconstruction
returns exactly the two original indices, and reading off an image recovers
all four parameters. Thus the cardinality is
$3\cdot3\cdot q\cdot q=9q^2$. The same bijection applies when $q=0$:
both sets are empty.
\end{proof}
